\documentclass[10pt,a4paper]{article}
\usepackage[margin=1.8cm]{geometry}
\usepackage[utf8]{inputenc}
\usepackage[T1]{fontenc}
\usepackage{enumitem}
\usepackage[hidelinks]{hyperref}
\usepackage{titlesec}
\pagestyle{empty}
\setlist[itemize]{leftmargin=1.2em,itemsep=1pt,topsep=2pt}
\titleformat{\section}{\large\bfseries}{}{0em}{}[\titlerule]
\titlespacing*{\section}{0pt}{10pt}{4pt}

\begin{document}
{\LARGE\bfseries Diogo Silva}\\[2pt]
Portugal \textbar{} \href{https://github.com/DaSilvaaa}{github.com/DaSilvaaa} \textbar{} \href{https://www.linkedin.com/in/diogo-silva-860185351}{linkedin.com/in/diogo-silva-860185351}

\section*{Profile}
Computer Engineering graduate (University of Aveiro) with a focus on software security, applied machine learning and systems programming, and author of an open monorepo with 18 tested projects in Python, Rust, C++, Java and C\#. Native European Portuguese speaker, English C1. Experience producing AI training data as a remote contractor.

\section*{Education}
\textbf{MSc in Sustainable Energy Systems}, University of Aveiro \hfill in progress\\
\textbf{BSc in Computer Engineering}, University of Aveiro \hfill 2022 -- 2026\\
Dissertation (18/20): prediction of vacancy formation energy in transition-metal-doped MXene surfaces (decision tree, XGBoost, variational autoencoder, SHAP).

\section*{Experience}
\textbf{Writer (Odin project)}, Mercor (remote, contract) \hfill Oct 2026 -- Nov 2026
\begin{itemize}
  \item Produced screen-recorded Python task demonstrations with step-by-step English narration for AI training data.
  \item Followed detailed task guidelines and incorporated quality review feedback.
\end{itemize}

\section*{Featured Projects}
\textbf{Portfolio\_Engineering (Polyglot Monorepo)} \hfill \href{https://github.com/DaSilvaaa/Portfolio_Engineering}{github.com/DaSilvaaa/Portfolio\_Engineering}
\begin{itemize}
  \item Monorepo with 18 standalone projects and a CI pipeline with 10 parallel GitHub Actions jobs, CodeQL, memory sanitizers (TSan/ASan/UBSan), gitleaks, Trivy and Dependabot.
  \item \textbf{Security \& Authentication:} hardened web API (Argon2id, JWT HS256 with a fixed algorithm, TOTP 2FA, password reset, anti-user-enumeration, rate limiting and a STRIDE threat model).
  \item \textbf{Forensic Tools \& Auditing:} entropy-based secret scanner using Shannon entropy (Rust), AST static analyzer with property-based tests (Hypothesis), TLS auditor (C\#) and AES-256-GCM vault with PBKDF2 (Java).
  \item \textbf{Systems \& Concurrency:} generic thread pool with \texttt{std::future} and thread-safe LRU cache validated against race conditions, bounds-checked PCAP packet parser (C++17).
  \item \textbf{Machine Learning \& MLOps:} leakage-free (\textit{no look-ahead}) hourly photovoltaic forecasting, network intrusion detection with Random Forest and Isolation Forest, leak-proof tabular pipeline with FastAPI inference and numerical implementations from scratch in NumPy.
\end{itemize}

\section*{Technical Skills}
\textbf{Languages:} Python, Rust, C++, Java, C\#, SQL, Bash\\
\textbf{Security \& DevSecOps:} STRIDE modelling, SAST, dependency auditing (pip-audit, cargo-audit), Docker\\
\textbf{Machine Learning:} scikit-learn, XGBoost, TensorFlow/Keras, SHAP, NumPy, pandas, time series, leak-proof cross-validation, FastAPI\\
\textbf{Tools:} Git, GitHub Actions, Docker, CMake, Maven, .NET SDK, pytest, Hypothesis

\section*{Certifications \& Courses}
Scrum Fundamentals Certified (SFC), SCRUMstudy, 2026 (ID 1199300). Courses: IBM SkillsBuild AI Fundamentals; Microsoft Learn AI concepts; AWS (Amazon Bedrock, ML \& AI fundamentals); Kaggle (Intermediate ML, Feature Engineering, ML Explainability); Anthropic (Claude API, MCP); Mercor (Introduction to Agents, Mastering Model Mechanics).
\end{document}
